\begin{definition}[Set Arithmetic Images]
\label{def:set-arithmetic-images}
Let $A,B\subseteq\mathbb{R}$ and let $c\in\mathbb{R}$.  The arithmetic
images of $A$ and $B$ under addition, subtraction, and multiplication are
defined by
\[
  A+B:=\{a+b:a\in A,\ b\in B\},
\]
\[
  A-B:=\{a-b:a\in A,\ b\in B\},
\]
and
\[
  AB:=\{ab:a\in A,\ b\in B\}.
\]
The scalar image of $A$ under multiplication by $c$ is
\[
  cA:=\{ca:a\in A\}.
\]
\end{definition}
\end{definitionbox}

\begin{remark*}[Standard quantified statement]
\[
x\in A+B\Longleftrightarrow(\exists a\in A)(\exists b\in B)(x=a+b),
\]
and similarly for difference and product.  For the scalar image,
\[
x\in cA\Longleftrightarrow(\exists a\in A)(x=ca).
\]
\end{remark*}

\begin{remark*}[Predicate reading]
Let
\[
D=A\times B,
\]
and define the arithmetic rules by
\[
\varphi_{+}(a,b)=a+b,\qquad
\varphi_{-}(a,b)=a-b,
\]
\[
\varphi_{\cdot}(a,b)=ab,\qquad
\varphi_c(a)=ca.
\]
Then
\[
\operatorname{ReplacementImage}(A+B,D,\varphi_{+}),
\]
\[
\operatorname{ReplacementImage}(A-B,D,\varphi_{-}),
\quad
\operatorname{ReplacementImage}(AB,D,\varphi_{\cdot}),
\]
\[
\operatorname{ReplacementImage}(cA,A,\varphi_c).
\]
\end{remark*}

\begin{remark*}[Negated quantified statement]
\[
x\notin A+B\Longleftrightarrow(\forall a\in A)(\forall b\in B)(x\neq a+b),
\]
with the analogous readings for the other arithmetic images.
\end{remark*}

\begin{remark*}[Negation predicate reading]
\[
C\ne AB
\quad\text{means}\quad
(\exists x)\bigl[x\in C\triangle AB\bigr],
\]
and similarly for sum, difference, product, and scalar image.
\end{remark*}

\begin{remark*}[Failure modes]
\begin{description}
  \item[Exposition.]
An arithmetic-image definition fails only if it omits a value produced by the
rule or includes a value not produced by the rule.


  \item[\exists-condition.]
  \[
x\notin A+B\Longleftrightarrow(\forall a\in A)(\forall b\in B)(x\neq a+b),
\]
\[
C\ne AB
\quad\text{means}\quad
(\exists x)\bigl[x\in C\triangle AB\bigr],
\]
\end{description}
\end{remark*}

\begin{remark*}[Failure modes]
\begin{description}
  \item[Exposition.]
Every membership claim decomposes into witnesses from the input sets and the
arithmetic operation that produced the image point.


  \item[\exists-condition.]
  \[
x\notin A+B\Longleftrightarrow(\forall a\in A)(\forall b\in B)(x\neq a+b),
\]
\[
C\ne AB
\quad\text{means}\quad
(\exists x)\bigl[x\in C\triangle AB\bigr],
\]
\end{description}
\end{remark*}

\begin{remark*}[Interpretation]
Set arithmetic turns pointwise arithmetic into images of sets. Membership in
one of these images is always witnessed by input elements from the original
sets and the corresponding arithmetic operation.
\end{remark*}

\begin{dependencies}
\begin{itemize}
  \item \hyperref[def:reals]{Real Numbers}
  \item \hyperref[def:subset]{Subset}
  \item \hyperref[def:cartesian-product]{Cartesian Product}
  \item \hyperref[def:image-set]{Image Set}
  \item \hyperref[def:addition-on-r]{Addition on $\mathbb{R}$}
  \item \hyperref[def:subtraction-on-r]{Subtraction on $\mathbb{R}$}
  \item \hyperref[def:multiplication-on-r]{Multiplication on $\mathbb{R}$}
\end{itemize}
\end{dependencies}

